\section{Continuing a lower barrier}\label{sec:barriers}

The local density--field comparison in \cref{prop:inverse-comparison} does not by itself
exclude a limiting field which is too small near the nucleus.  We now
carry a positive barrier from the nuclear scale to an arbitrary
intermediate scale.  The construction keeps the actual conditional
density inside the current radius and uses the Thomas--Fermi reaction
outside it.  Rare failures of the local comparison are retained as an
explicit, small source error.  This adapts the intermediate-subsolution
construction of \cite[\FIntermediate]{Foundation}, using the positive-price
comparison proved here to allow the growing offset $D_0=O(s^{-7})$.

Write $d=|y|$ and put
\begin{equation}\label{eq:barrier-function}
 f(t)=4\pi k(t_+)^{3/2},\qquad
 b_r(y)=B d^{-4}\left(1-\frac r{8d}\right)\quad(d\ge r),
\end{equation}
where $B>0$ will be fixed below.  Direct differentiation gives
\[
 \Delta b_r=Bd^{-6}\left(12-\frac{20r}{8d}\right)
       \ge \frac{19}{2}Bd^{-6}.
\]
Thus $\Delta b_r\ge f(b_r)$ whenever
$4\pi k\sqrt B\le19/2$.  If $R=2r$, then
\begin{equation}\label{eq:barrier-gap}
 b_r-b_R\ge\gamma R^{-4}\quad(R\le d\le1.1R),
 \qquad \gamma=\frac{B}{16(1.1)^5}>0.
\end{equation}

\begin{proposition}[Conditional barrier invariant]\label{prop:barrier-invariant}
There are fixed positive constants $\eps,B,C_{\mathrm{inv}}$ and $M>2$
with the following property.  For every sufficiently far system in
\cref{eq:scaling-regime}, use the observations and kernels of
\cref{eq:observation-scales,eq:master-kernel} with this $\eps$.
For every $0\le j\le J_*$ there exist $\mathcal F_j$-measurable
functions $u_j,p_j$ such that, with $r=r_j$,
\begin{equation}\label{eq:barrier-invariant}
 \begin{aligned}
 \Delta u_j&\ge-4\pi\nu
       +4\pi\ind_{\{d<r\}}\mu_j-4\pi p_j
       +\ind_{\{d\ge r\}}f(u_j),\\
 b_r\le u_j&\le C_{\mathrm{inv}}d^{-4}\quad(d\ge r),\\
 u_j&=b_r\quad(d>Mr).
 \end{aligned}
\end{equation}
The first inequality holds in distributions on $\R^3$.
For each fixed system, $u_j-V$ is continuous and locally Lipschitz,
with deterministic bounds on compact sets; $p_j$ is a nonnegative
bounded density supported in $\{d\le r_j\}$, also with a deterministic
bound.  Its expected total mass satisfies
\begin{equation}\label{eq:barrier-error-mass}
 \sup_{0\le j\le J_*}\E\int p_j
 \le C\bigl(r_0^{32}+s^{19/2-3w}\bigr)=o(1).
\end{equation}
The deterministic fixed-system regularity bounds need not be uniform
in $Z$; the constants in \cref{eq:barrier-error-mass} are uniform.
\end{proposition}

\begin{proof}
We use the weak maximum rule in \cite[\FMaximum]{Foundation}.
In its nonsingular form, if finitely many continuous locally $H^1$
functions satisfy
\[
 \Delta v_i\ge g(y)+F(y,v_i),
\]
where $g$ is locally bounded and $F$ is measurable in $y$, continuous
in its second variable, and locally bounded on bounded ranges, then
their maximum satisfies the same inequality with its own value in
the reaction.  The lemma also allows a common singular term
$-4\pi\nu$ when all branches have continuous locally $H^1$ offsets
from $V$ and the reaction vanishes near the origin.  These are the
precise forms used below.

\paragraph{Initialization.}
For sufficiently small fixed $\eps$, the initial conditional density
satisfies, outside an event of probability at most $Cr_0^{35}$,
\begin{equation}\label{eq:initial-potential}
 P_0(y):=K*(\ind_{\{d<r_0\}}\mu_0)(y)
       \le0.1\,Z/r_0
       \quad(r_0\le |y|\le1.1r_0).
\end{equation}
To prove this, fix $y$ in the indicated annulus and consider the
$\mathcal F_0$-event $A_y=\{P_0(y)>0.05Z/r_0\}$.  If its probability
is greater than $r_0^{40}$, \cref{lem:event-tilt} realizes its raw
conditional marginal as a normalized state in sector $n$ with
unshifted offset at most
\[
 D_0+Cr_0^{-2.02}\log^5(e/r_0^{40}).
\]
By \cref{eq:rough-deficit}, $D_0\le Cs^{-7}$.  Since $r_0/s\to0$,
the displayed offset is at most $Z^{7/3}$ eventually, with $\eps$
fixed.  The global density estimate of
\cite[\FGlobalDensity]{Foundation} applies: the state is normalized and
antisymmetric and $n\le3Z$, so its ordinary density $\rho_{A_y}$
obeys $\int\rho_{A_y}^{5/3}\le CZ^{7/3}$.

Only packet centers in $|x|<2r_0$ can contribute to the truncated
density defining $P_0$.  The potential of a truncated packet is at
most that of the full packet, and radial smearing bounds the latter
by $K(y-x)$.  The observation tower and H\"older's inequality give
\[
 \E[P_0(y)\mid A_y]
 \le\int_{|x|<2r_0}\frac{\rho_{A_y}(x)}{|y-x|}\,dx
 \le CZ^{7/5}r_0^{1/5}
 =C\eps^{6/5}Z/r_0.
\]
Choose $\eps$ so small that the last coefficient is less than $0.05$.
This contradicts the event definition.  Hence
$\Prob(A_y)\le r_0^{40}$.

The minimum packet width is $c_1r_0^{1+w}$.  Integrating
$|y-z|^{-2}$ against each bounded packet yields
\[
 \norm{\nabla P_0}_\infty\le CZr_0^{-2-2w}.
\]
A spatial net of spacing at most $r_0^2$ in the annulus has
$O(r_0^{-3})$ points.  Its union failure probability is
$O(r_0^{37})$, and its interpolation error divided by $Z/r_0$ is
$O(r_0^{1-2w})$.  This proves \cref{eq:initial-potential}, with the
weaker exponent stated there.  The good event is measurable because
$P_0$ is continuous in space.

Let $G_0$ be this good event and define
\[
 \mu^{\mathrm{in}}=\ind_{G_0}\ind_{\{d<r_0\}}\mu_0,
 \qquad p_0=\ind_{G_0^c}\ind_{\{d<r_0\}}\mu_0.
\]
On $d<2r_0$ consider the branch
\[
 v=V-K*\mu^{\mathrm{in}}-0.7Z/r_0
                +C_3(Z/r_0)^{3/2}d^2.
\]
Fix $C_3$ so that $6C_3\ge4\pi k\,2^{3/2}$ and decrease $\eps$
if necessary so that $4C_3\eps^{3/2}\le0.05$.  Since
$Zr_0^3=\eps^3$, the quadratic term is at most $0.05Z/r_0$ on
this ball.  On $r_0\le d<2r_0$ we have $v\le2Z/r_0$, and therefore
\[
 \Delta v=-4\pi\nu+4\pi\mu^{\mathrm{in}}
                   +6C_3(Z/r_0)^{3/2},
 \qquad 6C_3(Z/r_0)^{3/2}\ge f(v).
\]
On the lower overlap $r_0\le d<1.06r_0$,
\[
 v\ge(1/1.06-0.1-0.7)Z/r_0>0.14Z/r_0;
\]
on $G_0^c$ the same bound follows with the $0.1$ term omitted.
On $1.9r_0<d<2r_0$,
\[
 v\le(1/1.9-0.7+0.05)Z/r_0<0.
\]
Fix $B$ satisfying the subsolution condition above and
$0<B\le0.1\eps^3$.  Define $u_0$ on an open cover by
\[
 u_0=
 \begin{cases}
 v,&d<1.06r_0,\\
 \max(v,b_{r_0}),&r_0<d<2r_0,\\
 b_{r_0},&d>1.9r_0.
 \end{cases}
\]
The definitions agree on overlaps.  Near $d=r_0$ the branch $v$
alone is active and supplies both required source inequalities.
Elsewhere the weak maximum rule proves
\cref{eq:barrier-invariant}.  The exterior cap holds for
$C_{\mathrm{inv}}\ge\max(B,32\eps^3)$, and
\[
 \E\int p_0\le n\Prob(G_0^c)\le Cr_0^{32}.
\]

\paragraph{Choice of the remaining constants.}
Increase $C_{\mathrm{inv}}$ to exceed the universal cap
$C_{\mathrm{cap}}$ in \cref{prop:inverse-comparison}.
Choose
\[
 0<\lambda_2<\lambda_1<\gamma/2,
 \qquad M>2,\qquad C_{\mathrm{inv}}M^{-4}<\lambda_2.
\]
The constants $\lambda_1,\lambda_2$ are barrier shifts, distinct from
the electron price $\lambda=s^{-4}$.  Fix
$0<h_l<B/4$ and $h_h>C_{\mathrm{inv}}$, and then choose
$e_0,\xi>0$ so small that
\begin{equation}\label{eq:barrier-parameters}
 \begin{gathered}
 2e_0(2M)^4<B/4,\qquad
 16\xi+2e_0<\lambda_1-\lambda_2,\\
 \xi<\min(\lambda_2,h_l/16).
 \end{gathered}
\end{equation}
All these constants are now fixed.  The simultaneous comparison
\cref{prop:inverse-comparison} applies with exactly these parameters
along the regime under consideration.

\paragraph{One outward step.}
Suppose the invariant holds at $r=r_j$, where $j<J_*$, and set
$R=2r$ and $u=u_j$.  On the good event of
\cref{prop:inverse-comparison} for the band $r\le d\le4Mr$, use
\[
 v_1=u-\lambda_1R^{-4},\qquad
 v_2=H'_j-\lambda_2R^{-4}.
\]
The two branches have complementary roles.  On the gained annulus
$r\le d<R$, the old reaction must supply the actual density, so a retained
$v_1$ must satisfy $f(u)\ge4\pi\mu_j$.  Outside $R$, the field branch
already has Laplacian $4\pi\mu_j$, which must instead dominate $f(v_2)$.
The opposite inverse implications in \cref{eq:inverse-comparison},
together with the unequal shifts, give these two inequalities on the
neighborhoods used below.
On bad data omit $v_2$.  Define $\widetilde u$ by
\begin{equation}\label{eq:barrier-open-cover}
 \begin{array}{c|c}
 \text{region}&\widetilde u\\ \hline
 d<1.08R&\max(v_1,v_2)\\
 1.04R<d<2MR&\max(v_1,v_2,b_R)\\
 d>MR&b_R
 \end{array}
\end{equation}
with the same omission on bad data, and set
\begin{equation}\label{eq:barrier-added-error}
 \widetilde p=p_j+\ind_{\{\mathrm{bad}\}}\ind_{\{r\le d<R\}}\mu_j.
\end{equation}
\Cref{fig:barrier-step} shows the overlapping regions in
\cref{eq:barrier-open-cover} and the source-gain annulus $r\le d<R$.
On the lower overlap, \cref{eq:barrier-gap} gives
$v_1\ge b_r-\lambda_1R^{-4}>b_R$.  On the upper overlap, the caps
and $C_{\mathrm{inv}}M^{-4}<\lambda_2$ put both shifted branches
below $b_R$.  Thus the definitions agree.  They give
$b_R\le\widetilde u\le C_{\mathrm{inv}}d^{-4}$ on $d\ge R$ and
$\widetilde u=b_R$ on $d>MR$; for $R\le d<1.08R$ the lower bound
already follows from $v_1$.

\begin{figure}[tbp]
  \centering
  \begin{tikzpicture}[
    x=1cm,y=1cm,
    font=\fontsize{9}{11}\selectfont,
    line cap=round,
    line join=round
  ]
    \fill[orange!12] (1.6,0.65) rectangle (3.4,-4.65);
    \draw[orange!75!black,thick]
      (1.6,0.78) -- (1.6,0.96) -- (3.4,0.96) -- (3.4,0.78);
    \node[above,text=orange!60!black] at (2.5,0.99)
      {$r\le d<R$: source gain};

    \foreach \position in {1.6,3.4,4.9,6.4,9.3,11.8} {
      \draw[gray!50,densely dotted] (\position,-0.12) -- (\position,-4.65);
      \draw (\position,-0.07) -- (\position,0.07);
    }
    \draw[-{Latex[length=2mm]}] (0,0) -- (12.9,0);
    \draw (0,-0.07) -- (0,0.07);
    \node[above] at (0,0.1) {$0$};
    \node[above] at (1.6,0.1) {$r$};
    \node[above] at (3.4,0.1) {$R=2r$};
    \node[above] at (4.9,0.1) {$1.04R$};
    \node[above] at (6.4,0.1) {$1.08R$};
    \node[above] at (9.3,0.1) {$MR$};
    \node[above] at (11.8,0.1) {$2MR$};
    \node[anchor=north east] at (12.9,-0.15) {$d=|y|$};

    \filldraw[fill=blue!7,draw=blue!65!black,thick]
      (0,-0.9) rectangle (6.4,-1.8);
    \node[align=center] at (3.2,-1.35)
      {$\max(v_1,v_2)$\\[2pt]$d<1.08R$};

    \filldraw[fill=blue!7,draw=blue!65!black,thick]
      (4.9,-2.2) rectangle (11.8,-3.1);
    \node[align=center] at (8.35,-2.65)
      {$\max(v_1,v_2,b_R)$\\[2pt]$1.04R<d<2MR$};

    \fill[green!7] (9.3,-3.5) rectangle (12.9,-4.4);
    \draw[green!35!black,thick]
      (12.65,-3.5) -- (9.3,-3.5) -- (9.3,-4.4) -- (12.65,-4.4);
    \draw[green!35!black,thick,-{Latex[length=2mm]}]
      (12.4,-3.95) -- (12.9,-3.95);
    \node[align=center] at (10.9,-3.95)
      {$b_R$\\[2pt]$d>MR$};

    \node[anchor=north,align=center,text=black!70] at (6.45,-4.85)
      {Schematic radial regions, not to scale.\\
       The fields need not be radial.};
  \end{tikzpicture}
  \caption{The overlapping definitions in one continuation step from $r$ to
    $R=2r$. The shifted branches are
    $v_1=u_j-\lambda_1R^{-4}$ and
    $v_2=H'_j-\lambda_2R^{-4}$; the second branch is omitted on bad data.
    The source term changes on the shaded annulus $r\le d<R$.
    The barrier gap makes the first two definitions agree on
    $1.04R<d<1.08R$, and the caps make the last two agree on
    $MR<d<2MR$.}
  \label{fig:barrier-step}
\end{figure}


It remains to prove the weak inequality before averaging.  We do so
on open neighborhoods, to avoid differentiating across the sets
where the maximizing branch changes.  At a point in one of the
open regions in \cref{eq:barrier-open-cover}, retain the candidates
whose value is within $e_0R^{-4}$ of the maximum there.  By continuity
and finiteness, there is a neighborhood on which every discarded
candidate remains strictly below a fixed retained maximizer and
every retained candidate $v$ satisfies
\begin{equation}\label{eq:retained-branches}
 \widetilde u-v<2e_0R^{-4}.
\end{equation}
The maximum of the retained family is still $\widetilde u$ throughout
this neighborhood.  At the origin use the continuous offsets from
$V$ to make the same selection.

Every retained branch satisfies the common inequality
\begin{equation}\label{eq:common-branch-inequality}
 \Delta v\ge-4\pi\nu+4\pi\ind_{\{d<R\}}\mu_j
       -4\pi\widetilde p+\ind_{\{d\ge R\}}f(v)
\end{equation}
on this whole neighborhood.  For $v_1$, its old inequality suffices
on $d<r$, and on $d\ge R$ use $f(u)\ge f(v_1)$.
The only new source is on $r\le d<R$.  Bad data there are paid for
by \cref{eq:barrier-added-error}.  On good data $v_2$ is a candidate,
so \cref{eq:retained-branches} gives
\[
 H'_j\le u-(\lambda_1-\lambda_2-2e_0)R^{-4}.
\]
Moreover $h_l<7B/8\le d^4u\le C_{\mathrm{inv}}<h_h$.
If $4\pi\mu_j>f(u)$, the high implication of
\cref{eq:inverse-comparison}, at height $h=d^4u$, would give
\[
 H'_j\ge u-\xi d^{-4}\ge u-16\xi R^{-4},
\]
contrary to \cref{eq:barrier-parameters}.  Hence $4\pi\mu_j\le f(u)$,
which supplies the required new density.

For $v_2$ its exact identity
$\Delta v_2=-4\pi\nu+4\pi\mu_j$ suffices on $d<R$.
At a remaining point where it is retained, the lower splice or
the included barrier gives $\widetilde u\ge b_R$ and $d\le2MR$.
Consequently
\[
 d^4v_2\ge\frac{7B}{8}-2e_0(2M)^4>\frac{5B}{8}>h_l,
 \qquad d^4v_2\le C_{\mathrm{cap}}<h_h.
\]
If $4\pi\mu_j<f(v_2)$, the low implication of
\cref{eq:inverse-comparison} would imply
\[
 H'_j\le v_2+\xi d^{-4}\le v_2+\xi R^{-4}
                <v_2+\lambda_2R^{-4}=H'_j,
\]
a contradiction.  Thus $4\pi\mu_j\ge f(v_2)$ there.  The branch
$b_R$ occurs only outside $R$ and satisfies
\cref{eq:common-branch-inequality} by its subsolution property.

These comparisons concern zero-order densities almost everywhere
on the same open neighborhood, including when it crosses $d=R$.
The common reaction is $F(y,t)=\ind_{\{d\ge R\}}f(t)$.
It meets the weak maximum rule's hypotheses; near zero the common
$V$ singularity cancels and the reaction vanishes.  The rule therefore
proves \cref{eq:common-branch-inequality} for $\widetilde u$ locally.
A partition of unity proves it globally.

\paragraph{Forgetting one observation.}
Define
\[
 u_{j+1}=V+\E[\widetilde u-V\mid\mathcal F_{j+1}],
 \qquad
 p_{j+1}=\E[\widetilde p\mid\mathcal F_{j+1}].
\]
We justify the conditional integration of the weak inequality.
For a fixed system the positive minimum packet width gives
deterministic bounds for $\mu_j$, $K*\mu_j$, and its first
derivatives.  The initial truncated potential has the same bounds.
Thus the initial offsets have deterministic local Lipschitz bounds;
subtracting constants, taking finite maxima, and gluing on the fixed
open cover preserve such bounds.  Where $b_R$ is used, its offset
from $V$ is smooth and separated from the origin.  The errors remain
bounded densities supported in deterministic balls.  Conditional
integration preserves these bounds.  All fields and flags are jointly
measurable; the spatial supremum defining the initial flag can be
tested on a countable dense set.  Fubini's theorem therefore permits
conditional integration against each smooth compactly supported
test.  A countable dense family of nonnegative tests, followed by
local approximation, supplies simultaneous almost-sure weak
inequalities for all tests.

The conditional tower gives $\mu_{j+1}$ as the new inner density,
and convexity gives
\[
 \E[f(\widetilde u)\mid\mathcal F_{j+1}]
 \ge f\bigl(V+\E[\widetilde u-V\mid\mathcal F_{j+1}]\bigr).
\]
This proves the first line of \cref{eq:barrier-invariant} at $R$.
The deterministic lower and upper bounds, the equality beyond
$MR$, and the error support also pass to the conditional average.

\paragraph{Total error.}
For $r\le |y|<2r$, every master packet contributing to $\mu_j(y)$
has center in a fixed annulus comparable to $r$ and width comparable
to $r^{1+w}$.  Conditional Jensen, the packet bound, and
\cref{eq:annular-count,eq:rough-deficit} imply
\[
 \norm{\mu_j(y)}_{L^2(\Prob)}
 \le Cr^{-3-3w}\bigl(r^{-3}+\sqrt{D_0r}\bigr)
 \le Cr^{-6-3w},
\]
since $r\le s$ and $D_0\le Cs^{-7}$.  The bad band event has
probability at most $Cr^{25}$.  Cauchy--Schwarz in the data and
integration over the gain annulus therefore give
\[
 \E\int\ind_{\{\mathrm{bad}\}}\ind_{\{r\le d<R\}}\mu_j
 \le Cr^{25/2+3-6-3w}=Cr^{19/2-3w}.
\]
The sum over the geometric sequence of scales is at most
$Cs^{19/2-3w}$.  Conditional averaging preserves expected total mass.
Together with the initial error estimate this proves
\cref{eq:barrier-error-mass} and completes the induction.
\end{proof}
